\documentclass[exercice]{thedocument}%
\usepackage{texmaths-tikz}%
\startdatakeys
\source{}
\cycle{}
\competencesDuSocle{}
\connaissancesRequises{}
\competencesTravaillees{}
\enddatakeys
\begin{document}%
Enzo affirme : \guillemets{La droite $(d)$ est un axe de
symétrie de cette figure tracée à main levée.}\par\noindent
Cette affirmation est-elle vraie ou fausse ?
\begin{center}
    \begin{tikzpicture}[scale=0.9]
        \coordinate[label=above left:{\<S1;}](S1)at(-2,2);
        \coordinate[label=above right:{\<S2;}](S2)at(2,2);
        \coordinate[label=below right:{\<S3;}](S3)at(2,-1);
        \coordinate[label=below left:{\<S4;}](S4)at(-2,-1);
        \coordinate(I)at(0,0);
        \coordinate(J)at(0,2);
        \coordinate[label=right:$\sf(d)$](d0)at(0,3);
        \coordinate(d1)at(0,-1.5);
        \tkzMarkRightAngle[color=red](I,J,S1)
        \draw(S1)--(S2)--(S3)--(I)--(S4)--cycle;
        \draw[red](d0)--(d1);
        \path(S1)--(J)node[pos=0.5]{\color{red}\bfseries/};
        \path(S2)--(J)node[pos=0.5]{\color{red}\bfseries/};
        \path(S4)--(S1)node[pos=0.5, above, sloped]{5 cm};
        \path(S4)--(I)node[pos=0.5, below, sloped]{3 cm};
        \path(I)--(S3)node[pos=0.5, below, sloped]{3 cm};
        \path(S3)--(S2)node[pos=0.5, below, sloped]{4 cm};
    \end{tikzpicture}
\end{center}
\end{document}
